\documentclass[12pt,a4paper]{article}

\usepackage[T2A]{fontenc}
\usepackage[utf8]{inputenc}
\usepackage[russian]{babel}

\usepackage{amsmath, amssymb, amsthm}
\usepackage{mathtools}
\usepackage{stmaryrd} % для \llbracket и \rrbracket

\usepackage{enumitem}
\usepackage{array}

\usepackage{xcolor}
\usepackage{tikz}
\usetikzlibrary{
    arrows.meta,
    automata,
    positioning,
    shapes.geometric % для diamond и aspect
}

\usepackage{listings}
\usepackage{hyperref}

\usetikzlibrary{arrows.meta, automata, positioning}

\newtheorem{definition}{Определение}
\newtheorem{theorem}{Теорема}
\newtheorem{lemma}{Лемма}
\newtheorem{example}{Пример}
\newtheorem{remark}{Замечание}

\lstdefinestyle{pseudo}{
    basicstyle=\ttfamily\small,
    keywordstyle=\color{blue},
    commentstyle=\color{gray},
    numbers=left,
    numberstyle=\tiny,
    frame=single,
    breaklines=true,
    tabsize=2,
    captionpos=b,
    columns=fullflexible,
    keepspaces=true,
    inputencoding=utf8,
    extendedchars=true,
    literate=
        {А}{{\CYRA}}1
        {Б}{{\CYRB}}1
        {В}{{\CYRV}}1
        {Г}{{\CYRG}}1
        {Д}{{\CYRD}}1
        {Е}{{\CYRE}}1
        {Ё}{{\CYRYO}}1
        {Ж}{{\CYRZH}}1
        {З}{{\CYRZ}}1
        {И}{{\CYRI}}1
        {Й}{{\CYRISHRT}}1
        {К}{{\CYRK}}1
        {Л}{{\CYRL}}1
        {М}{{\CYRM}}1
        {Н}{{\CYRN}}1
        {О}{{\CYRO}}1
        {П}{{\CYRP}}1
        {Р}{{\CYRR}}1
        {С}{{\CYRS}}1
        {Т}{{\CYRT}}1
        {У}{{\CYRU}}1
        {Ф}{{\CYRF}}1
        {Х}{{\CYRH}}1
        {Ц}{{\CYRC}}1
        {Ч}{{\CYRCH}}1
        {Ш}{{\CYRSH}}1
        {Щ}{{\CYRSHCH}}1
        {Ъ}{{\CYRHRDSN}}1
        {Ы}{{\CYRERY}}1
        {Ь}{{\CYRSFTSN}}1
        {Э}{{\CYREREV}}1
        {Ю}{{\CYRYU}}1
        {Я}{{\CYRYA}}1
        {а}{{\cyra}}1
        {б}{{\cyrb}}1
        {в}{{\cyrv}}1
        {г}{{\cyrg}}1
        {д}{{\cyrd}}1
        {е}{{\cyre}}1
        {ё}{{\cyryo}}1
        {ж}{{\cyrzh}}1
        {з}{{\cyrz}}1
        {и}{{\cyri}}1
        {й}{{\cyrishrt}}1
        {к}{{\cyrk}}1
        {л}{{\cyrl}}1
        {м}{{\cyrm}}1
        {н}{{\cyrn}}1
        {о}{{\cyro}}1
        {п}{{\cyrp}}1
        {р}{{\cyrr}}1
        {с}{{\cyrs}}1
        {т}{{\cyrt}}1
        {у}{{\cyru}}1
        {ф}{{\cyrf}}1
        {х}{{\cyrh}}1
        {ц}{{\cyrc}}1
        {ч}{{\cyrch}}1
        {ш}{{\cyrsh}}1
        {щ}{{\cyrshch}}1
        {ъ}{{\cyrhrdsn}}1
        {ы}{{\cyrery}}1
        {ь}{{\cyrsftsn}}1
        {э}{{\cyrerev}}1
        {ю}{{\cyryu}}1
        {я}{{\cyrya}}1
}

\title{Билет №42 \\ \small Методы спецификации программ. Схемное, структурное, визуальное, автоматное программирование. Отличие управляющих автоматов от абстрактных. (ИТМО) \\ Методы спецификации программ. Методы проверки спецификации. (СпБПУ)}
\author{Сериков Владислав Александрович}
\date{13 мая 2026}

\begin{document}

\maketitle
\tableofcontents
\newpage

\section{Понятие спецификации программы}

\begin{definition}
\textbf{Спецификация программы} --- это описание требований к программе, её интерфейсам, поведению, ограничениям и ожидаемым свойствам, не обязательно совпадающее с текстом реализации.
\end{definition}

Спецификация отвечает на вопрос: \emph{что должна делать программа}, тогда как реализация отвечает на вопрос: \emph{как именно программа это делает}.

Обычно спецификация включает:
\begin{enumerate}
    \item назначение программы;
    \item входные и выходные данные;
    \item предусловия и постусловия;
    \item инварианты;
    \item ограничения на время, память, безопасность, надёжность;
    \item реакции на ошибки;
    \item допустимые последовательности взаимодействий с внешней средой.
\end{enumerate}

\begin{definition}
\textbf{Формальная спецификация} --- это спецификация, записанная на языке с математически определённым синтаксисом и семантикой.
\end{definition}

Примеры формальных средств спецификации:
\begin{itemize}
    \item логика предикатов;
    \item алгебраические спецификации;
    \item автоматы;
    \item сети Петри;
    \item временные логики LTL и CTL;
    \item спецификационные языки Z, VDM, B, TLA+;
    \item контракты вида \texttt{requires/ensures/invariant}.
\end{itemize}

\section{Классификация методов спецификации программ}

Методы спецификации можно условно разделить на несколько групп.

\subsection{Неформальные спецификации}

Неформальная спецификация записывается на естественном языке.

\textbf{Достоинства:}
\begin{itemize}
    \item понятна широкому кругу участников проекта;
    \item не требует специальной математической подготовки;
    \item удобна на ранних стадиях анализа требований.
\end{itemize}

\textbf{Недостатки:}
\begin{itemize}
    \item неоднозначность;
    \item неполнота;
    \item противоречивость;
    \item трудность автоматической проверки.
\end{itemize}

\subsection{Полуформальные спецификации}

Полуформальные спецификации используют графические или табличные нотации, семантика которых частично формализована.

Примеры:
\begin{itemize}
    \item UML-диаграммы;
    \item блок-схемы;
    \item диаграммы состояний;
    \item таблицы решений;
    \item диаграммы потоков данных;
    \item BPMN-модели.
\end{itemize}

Такие спецификации удобны для общения с заказчиком и проектирования архитектуры, но без строгой семантики они не всегда пригодны для доказательства корректности.

\subsection{Формальные спецификации}

Формальные спецификации задаются математически.

Примеры:
\begin{itemize}
    \item спецификация функции через предусловие и постусловие;
    \item автоматная модель протокола;
    \item описание системы во временной логике;
    \item аксиоматическое описание абстрактного типа данных;
    \item спецификация переходной системы.
\end{itemize}

\begin{example}
Спецификация функции деления:
\[
    \texttt{div}(a,b)
\]
может быть задана контрактом:
\[
    b \neq 0
\]
--- предусловие,
\[
    result = \frac{a}{b}
\]
--- постусловие.
\end{example}

\section{Спецификация через предусловия, постусловия и инварианты}

Одним из наиболее распространённых способов формальной спецификации является контрактная спецификация.

\begin{definition}
\textbf{Предусловие} --- это условие, которое должно быть истинно перед выполнением операции.
\end{definition}

\begin{definition}
\textbf{Постусловие} --- это условие, которое должно быть истинно после завершения операции, если перед её выполнением было истинно предусловие.
\end{definition}

\begin{definition}
\textbf{Инвариант} --- это условие, сохраняющее истинность в определённых точках выполнения программы.
\end{definition}

Для фрагмента программы $S$ корректность часто записывают в форме тройки Хоара:
\[
    \{P\}\ S\ \{Q\},
\]
где:
\begin{itemize}
    \item $P$ --- предусловие;
    \item $S$ --- программа или оператор;
    \item $Q$ --- постусловие.
\end{itemize}

Смысл записи:
если перед выполнением $S$ истинно $P$ и выполнение $S$ завершается, то после выполнения истинно $Q$.

\section{Частичная и полная корректность}

\begin{definition}
Программа $S$ называется \textbf{частично корректной} относительно предусловия $P$ и постусловия $Q$, если из истинности $P$ перед запуском и завершения $S$ следует истинность $Q$ после завершения.
\end{definition}

\[
    P \land S\downarrow \Rightarrow Q.
\]

Здесь $S\downarrow$ означает, что программа завершилась.

\begin{definition}
Программа $S$ называется \textbf{полностью корректной} относительно $P$ и $Q$, если она частично корректна и обязательно завершается при любом входе, удовлетворяющем $P$.
\end{definition}

Иначе говоря:
\[
    P \Rightarrow (S\downarrow \land Q).
\]

\section{Аксиоматический метод спецификации и логика Хоара}

Логика Хоара задаёт правила рассуждения о корректности программ.

\subsection{Правило присваивания}

\begin{theorem}[Правило присваивания в логике Хоара]
Пусть $Q$ --- утверждение о состоянии программы после присваивания, а $Q[E/x]$ обозначает формулу, полученную из $Q$ заменой всех свободных вхождений переменной $x$ на выражение $E$.
Тогда справедлива тройка Хоара:
\[
    \{Q[E/x]\}\ x := E\ \{Q\}.
\]
\end{theorem}

\begin{proof}
Рассмотрим произвольное состояние $\sigma$ перед выполнением оператора $x := E$.
Пусть в состоянии $\sigma$ истинна формула $Q[E/x]$.

После выполнения присваивания получается новое состояние $\sigma'$, которое отличается от $\sigma$ только значением переменной $x$:
\[
    \sigma'(x) = \llbracket E \rrbracket_\sigma,
\]
а для любой другой переменной $y \neq x$:
\[
    \sigma'(y) = \sigma(y).
\]

Формула $Q[E/x]$ в старом состоянии $\sigma$ утверждает ровно то же, что формула $Q$ в новом состоянии $\sigma'$, поскольку все места, где после присваивания используется новое значение $x$, до присваивания соответствуют значению выражения $E$.

Следовательно, если до выполнения оператора было истинно $Q[E/x]$, то после выполнения оператора истинно $Q$.
Значит:
\[
    \{Q[E/x]\}\ x := E\ \{Q\}.
\]
\end{proof}

\begin{example}
Пусть требуется доказать:
\[
    \{x + 1 > 0\}\ y := x + 1\ \{y > 0\}.
\]

Здесь:
\[
    Q \equiv y > 0,
\]
\[
    E \equiv x + 1.
\]

Заменим $y$ на $x+1$:
\[
    Q[x+1/y] \equiv x + 1 > 0.
\]

Получаем корректную тройку:
\[
    \{x + 1 > 0\}\ y := x + 1\ \{y > 0\}.
\]
\end{example}

\subsection{Правило последовательной композиции}

Если:
\[
    \{P\}\ S_1\ \{R\}
\]
и
\[
    \{R\}\ S_2\ \{Q\},
\]
то:
\[
    \{P\}\ S_1;S_2\ \{Q\}.
\]

\subsection{Правило ветвления}

Если:
\[
    \{P \land B\}\ S_1\ \{Q\}
\]
и
\[
    \{P \land \neg B\}\ S_2\ \{Q\},
\]
то:
\[
    \{P\}\ \textbf{if } B \textbf{ then } S_1 \textbf{ else } S_2\ \{Q\}.
\]

\subsection{Правило цикла}

Если:
\[
    \{I \land B\}\ S\ \{I\},
\]
то:
\[
    \{I\}\ \textbf{while } B \textbf{ do } S\ \{I \land \neg B\}.
\]

Здесь $I$ --- инвариант цикла.

\begin{example}
Рассмотрим программу вычисления суммы чисел от $1$ до $n$:

\begin{lstlisting}[style=pseudo]
i := 0;
s := 0;
while i < n do
    i := i + 1;
    s := s + i
end
\end{lstlisting}

Спецификация:
\[
    n \geq 0
\]
--- предусловие,
\[
    s = \frac{n(n+1)}{2}
\]
--- постусловие.

Инвариант:
\[
    I \equiv 0 \leq i \leq n \land s = \frac{i(i+1)}{2}.
\]

Проверка:
\begin{enumerate}
    \item До цикла: $i=0$, $s=0$, следовательно
    \[
        s = 0 = \frac{0(0+1)}{2}.
    \]
    Инвариант истинен.

    \item Если $I$ истинен и $i<n$, то после $i := i+1$ и $s := s+i$ получаем:
    \[
        s' = \frac{i(i+1)}{2} + (i+1)
            = \frac{(i+1)(i+2)}{2}.
    \]
    Значит, инвариант сохраняется.

    \item После выхода из цикла:
    \[
        I \land \neg(i<n).
    \]
    Из $0 \leq i \leq n$ и $i \geq n$ следует $i=n$.
    Тогда:
    \[
        s = \frac{n(n+1)}{2}.
    \]
\end{enumerate}
\end{example}

\section{Схемное программирование}

\begin{definition}
\textbf{Схемное программирование} --- это подход к построению программ, при котором программа задаётся в виде схемы, состоящей из операторных вершин, условий, переходов и управляющих связей.
\end{definition}

Исторически схемное программирование связано с блок-схемами, граф-схемами алгоритмов и схемами программ.

\subsection{Граф-схема алгоритма}

\begin{definition}
\textbf{Граф-схема алгоритма} --- это ориентированный граф, вершины которого соответствуют действиям или проверкам условий, а дуги задают передачу управления.
\end{definition}

Обычно используются следующие элементы:
\begin{itemize}
    \item начало и конец;
    \item процесс;
    \item условие;
    \item ввод/вывод;
    \item переход управления.
\end{itemize}

\begin{center}
\begin{tikzpicture}[node distance=1.5cm, >=Latex]
    \node[draw, rounded corners] (start) {Начало};
    \node[draw, below of=start] (input) {$n$};
    \node[draw, below of=input] (init) {$i:=1,\ s:=0$};
    \node[draw, diamond, aspect=2, below of=init, yshift=-0.3cm] (cond) {$i \leq n$?};
    \node[draw, right of=cond, xshift=3cm] (body) {$s:=s+i;\ i:=i+1$};
    \node[draw, below of=cond, yshift=-0.5cm] (out) {Вывести $s$};
    \node[draw, rounded corners, below of=out] (end) {Конец};

    \draw[->] (start) -- (input);
    \draw[->] (input) -- (init);
    \draw[->] (init) -- (cond);
    \draw[->] (cond) -- node[left] {нет} (out);
    \draw[->] (out) -- (end);
    \draw[->] (cond) -- node[above] {да} (body);
    \draw[->] (body) |- (cond);
\end{tikzpicture}
\end{center}

\subsection{Достоинства схемного подхода}

\begin{itemize}
    \item наглядность;
    \item удобство при объяснении алгоритмов;
    \item явное отображение управляющего потока;
    \item возможность анализа достижимости фрагментов схемы.
\end{itemize}

\subsection{Недостатки схемного подхода}

\begin{itemize}
    \item громоздкость для больших программ;
    \item склонность к неструктурированным переходам;
    \item сложность сопровождения;
    \item возможность возникновения ``спагетти-логики''.
\end{itemize}

\section{Структурное программирование}

\begin{definition}
\textbf{Структурное программирование} --- это методология разработки программ, при которой программа строится из ограниченного набора управляющих конструкций: последовательности, ветвления и цикла.
\end{definition}

Основная идея структурного программирования состоит в отказе от произвольных переходов \texttt{goto} и представлении программы как иерархии вложенных блоков.

\subsection{Базовые управляющие конструкции}

\subsubsection{Последовательность}

\[
    S_1; S_2
\]

Сначала выполняется $S_1$, затем $S_2$.

\subsubsection{Ветвление}

\[
    \textbf{if } B \textbf{ then } S_1 \textbf{ else } S_2.
\]

Если условие $B$ истинно, выполняется $S_1$, иначе --- $S_2$.

\subsubsection{Цикл}

\[
    \textbf{while } B \textbf{ do } S.
\]

Пока условие $B$ истинно, выполняется тело цикла $S$.

\section{Теорема Бёма--Якопини}

\begin{theorem}[Бём--Якопини]
Всякая вычислимая функция, реализуемая блок-схемой с одним входом и одним выходом, может быть реализована программой, использующей только три управляющие конструкции:
\begin{enumerate}
    \item последовательное выполнение;
    \item ветвление;
    \item цикл.
\end{enumerate}
\end{theorem}

\begin{proof}
Докажем конструктивно, что произвольную блок-схему с конечным числом вершин можно смоделировать структурной программой.

Пусть дана блок-схема с вершинами:
\[
    v_1, v_2, \ldots, v_n.
\]

Каждая вершина соответствует некоторому элементарному действию или условной проверке.
Введём специальную переменную управления:
\[
    pc,
\]
которая хранит номер текущей вершины, то есть аналог счётчика команд.

Начальному состоянию соответствует:
\[
    pc := 1.
\]

Конечной вершине поставим в соответствие значение:
\[
    pc := 0.
\]

Теперь построим программу следующего вида:

\begin{lstlisting}[style=pseudo]
pc := 1;
while pc != 0 do
    if pc == 1 then
        выполнить действие вершины v1;
        pc := номер следующей вершины
    else if pc == 2 then
        выполнить действие вершины v2;
        pc := номер следующей вершины
    ...
    else if pc == n then
        выполнить действие вершины vn;
        pc := номер следующей вершины
    end
end
\end{lstlisting}

Если вершина является операторной, то после выполнения соответствующего действия переменной $pc$ присваивается номер следующей вершины.

Если вершина является условной, то её можно записать с помощью ветвления:

\begin{lstlisting}[style=pseudo]
if B then
    pc := k
else
    pc := m
end
\end{lstlisting}

где $k$ и $m$ --- номера вершин, в которые ведут дуги из условной вершины.

Таким образом, вся блок-схема моделируется с помощью:
\begin{itemize}
    \item одной внешней конструкции \texttt{while};
    \item последовательного выполнения операторов;
    \item вложенной цепочки условных операторов \texttt{if}.
\end{itemize}

Поскольку каждый шаг построенной программы точно имитирует один шаг исходной блок-схемы, последовательность состояний данных в обеих моделях совпадает.
Следовательно, если исходная блок-схема вычисляет некоторую функцию, то построенная структурная программа вычисляет ту же функцию.

Теорема доказана.
\end{proof}

\begin{remark}
Теорема Бёма--Якопини показывает принципиальную достаточность трёх конструкций, но не утверждает, что механическое преобразование любой блок-схемы в такую форму всегда даёт удобную для чтения программу.
\end{remark}

\section{Визуальное программирование}

\begin{definition}
\textbf{Визуальное программирование} --- это подход, при котором программа или её существенная часть создаётся с помощью графических элементов: блоков, диаграмм, схем, соединений, состояний и переходов.
\end{definition}

Примеры визуальных средств:
\begin{itemize}
    \item блок-схемы;
    \item UML State Machine;
    \item UML Activity Diagram;
    \item Scratch;
    \item Simulink;
    \item LabVIEW;
    \item диаграммы состояний;
    \item языки для программируемых логических контроллеров, например SFC.
\end{itemize}

\subsection{Особенности визуального программирования}

\begin{itemize}
    \item программа представляется в виде графа;
    \item элементы графа имеют определённую семантику;
    \item связи задают поток данных или поток управления;
    \item часто возможна автоматическая генерация кода;
    \item удобно моделировать реактивные и событийные системы.
\end{itemize}

\subsection{Преимущества}

\begin{itemize}
    \item высокая наглядность;
    \item удобство для предметных специалистов;
    \item возможность визуальной отладки;
    \item хорошая применимость к автоматным моделям.
\end{itemize}

\subsection{Недостатки}

\begin{itemize}
    \item большие схемы быстро становятся трудно читаемыми;
    \item текстовые изменения иногда выполнять проще;
    \item затруднено версионирование;
    \item при нестрогой семантике возможны неоднозначности.
\end{itemize}

\section{Автоматное программирование}

\begin{definition}
\textbf{Автоматное программирование} --- это метод разработки программ, при котором поведение программы или её части описывается в терминах состояний, событий, переходов и действий.
\end{definition}

Автоматный подход особенно удобен для:
\begin{itemize}
    \item протоколов;
    \item пользовательских интерфейсов;
    \item встроенных систем;
    \item систем управления;
    \item лексических анализаторов;
    \item реактивных и событийных программ.
\end{itemize}

\subsection{Конечный автомат}

\begin{definition}
\textbf{Детерминированный конечный автомат} задаётся пятёркой:
\[
    A = (Q, \Sigma, \delta, q_0, F),
\]
где:
\begin{itemize}
    \item $Q$ --- конечное множество состояний;
    \item $\Sigma$ --- входной алфавит;
    \item $\delta: Q \times \Sigma \to Q$ --- функция переходов;
    \item $q_0 \in Q$ --- начальное состояние;
    \item $F \subseteq Q$ --- множество допускающих состояний.
\end{itemize}
\end{definition}

Такой автомат используется, например, для распознавания языков.

\subsection{Автомат Мили}

\begin{definition}
\textbf{Автомат Мили} задаётся шестёркой:
\[
    A = (Q, X, Y, \delta, \lambda, q_0),
\]
где:
\begin{itemize}
    \item $Q$ --- множество состояний;
    \item $X$ --- множество входных символов;
    \item $Y$ --- множество выходных символов;
    \item $\delta: Q \times X \to Q$ --- функция переходов;
    \item $\lambda: Q \times X \to Y$ --- функция выходов;
    \item $q_0$ --- начальное состояние.
\end{itemize}
\end{definition}

В автомате Мили выход зависит от текущего состояния и входного воздействия.

\subsection{Автомат Мура}

\begin{definition}
\textbf{Автомат Мура} задаётся шестёркой:
\[
    A = (Q, X, Y, \delta, \lambda, q_0),
\]
где:
\[
    \delta: Q \times X \to Q,
\]
а функция выходов имеет вид:
\[
    \lambda: Q \to Y.
\]
\end{definition}

В автомате Мура выход зависит только от текущего состояния.

\subsection{Пример автоматной спецификации}

Рассмотрим турникет метро.

Состояния:
\[
    Q = \{\text{Locked}, \text{Unlocked}\}.
\]

События:
\[
    X = \{\text{coin}, \text{push}\}.
\]

Переходы:

\[
\begin{array}{c|c|c|c}
\text{Состояние} & \text{Событие} & \text{Новое состояние} & \text{Действие} \\
\hline
\text{Locked} & \text{coin} & \text{Unlocked} & \text{разблокировать} \\
\text{Locked} & \text{push} & \text{Locked} & \text{сигнал отказа} \\
\text{Unlocked} & \text{push} & \text{Locked} & \text{заблокировать} \\
\text{Unlocked} & \text{coin} & \text{Unlocked} & \text{вернуть монету}
\end{array}
\]

Граф автомата:

\begin{center}
\begin{tikzpicture}[>=Latex, node distance=4cm, on grid, auto]
    \node[state, initial] (locked) {Locked};
    \node[state, right of=locked] (unlocked) {Unlocked};

    \path[->]
        (locked) edge[bend left] node {coin / unlock} (unlocked)
        (unlocked) edge[bend left] node {push / lock} (locked)
        (locked) edge[loop above] node {push / alarm} ()
        (unlocked) edge[loop above] node {coin / return} ();
\end{tikzpicture}
\end{center}

\subsection{Алгоритм реализации конечного автомата через таблицу переходов}

\textbf{Вход:}
\begin{itemize}
    \item текущее состояние $q$;
    \item событие $x$;
    \item таблица переходов.
\end{itemize}

\textbf{Выход:}
\begin{itemize}
    \item новое состояние;
    \item выполненное действие.
\end{itemize}

\begin{lstlisting}[style=pseudo]
state := initial_state

while программа работает do
    event := получить_событие()

    transition := таблица[state, event]

    if transition существует then
        выполнить transition.action
        state := transition.next_state
    else
        обработать_ошибку
    end
end
\end{lstlisting}

Минимальный пример на псевдокоде:

\begin{lstlisting}[style=pseudo]
state := Locked

on event coin:
    if state == Locked then
        unlock()
        state := Unlocked
    else if state == Unlocked then
        return_coin()
        state := Unlocked

on event push:
    if state == Locked then
        alarm()
        state := Locked
    else if state == Unlocked then
        lock()
        state := Locked
\end{lstlisting}

\section{Отличие управляющих автоматов от абстрактных автоматов}

\subsection{Абстрактный автомат}

\begin{definition}
\textbf{Абстрактный автомат} --- это математическая модель дискретного устройства, преобразующего входные последовательности в выходные последовательности и изменяющего состояние по заданным функциям переходов и выходов.
\end{definition}

Абстрактный автомат обычно задаётся на уровне множеств:
\[
    A = (Q, X, Y, \delta, \lambda, q_0).
\]

Он описывает:
\begin{itemize}
    \item множество состояний;
    \item входной алфавит;
    \item выходной алфавит;
    \item функцию переходов;
    \item функцию выходов.
\end{itemize}

При этом не уточняется, как именно автомат реализован аппаратно или программно.

\subsection{Управляющий автомат}

\begin{definition}
\textbf{Управляющий автомат} --- это автомат, предназначенный для управления некоторым операционным устройством, объектом управления или вычислительным процессом.
\end{definition}

Управляющий автомат:
\begin{itemize}
    \item принимает сигналы условий от операционной части;
    \item формирует управляющие воздействия;
    \item определяет последовательность выполнения операций;
    \item координирует работу системы.
\end{itemize}

Часто систему разделяют на две части:
\[
    \text{система} = \text{управляющий автомат} + \text{операционный автомат}.
\]

\begin{definition}
\textbf{Операционный автомат} --- это часть системы, выполняющая преобразование данных под управлением управляющего автомата.
\end{definition}

\subsection{Сравнение абстрактного и управляющего автомата}

\[
\begin{array}{|p{5cm}|p{5cm}|p{5cm}|}
\hline
\textbf{Критерий} & \textbf{Абстрактный автомат} & \textbf{Управляющий автомат} \\
\hline
Назначение & Математическое описание поведения & Управление действиями операционной части или объекта \\
\hline
Входы & Символы входного алфавита & События, признаки, флаги, условия от объекта управления \\
\hline
Выходы & Символы выходного алфавита & Управляющие сигналы, команды, вызовы действий \\
\hline
Уровень описания & Теоретико-автоматный & Архитектурный или инженерный \\
\hline
Связь с данными & Данные часто абстрагированы & Данные обрабатываются операционным автоматом \\
\hline
Роль & Распознавание или преобразование последовательностей & Организация порядка выполнения операций \\
\hline
\end{array}
\]

\subsection{Пример разделения на управляющий и операционный автоматы}

Задача: вычислить модуль числа $x$.

Операционная часть умеет:
\begin{itemize}
    \item проверять условие $x < 0$;
    \item выполнять операцию $x := -x$;
    \item выдавать результат.
\end{itemize}

Управляющий автомат задаёт порядок:

\[
\begin{array}{c|c|c|c}
\text{Состояние} & \text{Условие} & \text{Команда} & \text{Следующее состояние} \\
\hline
q_0 & x < 0 & x := -x & q_1 \\
q_0 & x \geq 0 & \text{ничего} & q_1 \\
q_1 & - & \text{выдать } x & q_f
\end{array}
\]

Здесь управляющий автомат не вычисляет сам модуль как арифметическое устройство, а определяет, какую операцию выполнить и когда.

\section{Методы проверки спецификаций}

Проверка спецификации направлена на обнаружение ошибок до реализации или до эксплуатации системы.

Основные свойства качественной спецификации:
\begin{enumerate}
    \item полнота;
    \item непротиворечивость;
    \item однозначность;
    \item реализуемость;
    \item проверяемость;
    \item трассируемость к требованиям;
    \item минимальность, то есть отсутствие избыточных требований.
\end{enumerate}

\subsection{Инспекция и рецензирование}

\begin{definition}
\textbf{Инспекция спецификации} --- это систематическая экспертная проверка документа спецификации с целью обнаружить ошибки, пропуски и противоречия.
\end{definition}

Проверяются:
\begin{itemize}
    \item неоднозначные формулировки;
    \item противоречивые требования;
    \item неполные сценарии;
    \item неописанные исключения;
    \item несогласованность терминов.
\end{itemize}

Типовой алгоритм инспекции:

\begin{enumerate}
    \item определить участников проверки;
    \item раздать спецификацию заранее;
    \item каждый участник независимо ищет дефекты;
    \item провести общее обсуждение;
    \item зафиксировать найденные дефекты;
    \item исправить спецификацию;
    \item провести повторную проверку.
\end{enumerate}

\subsection{Проверка полноты}

Спецификация считается полной, если она описывает все существенные случаи поведения.

Для функции:
\[
    f: D \to R
\]
необходимо описать результат для всех допустимых входов:
\[
    \forall x \in D \ \exists y \in R: y = f(x).
\]

\begin{example}
Неполная спецификация:
\begin{quote}
Если пользователь ввёл правильный пароль, система открывает доступ.
\end{quote}

Не описано, что происходит при неправильном пароле, пустом пароле, заблокированной учётной записи, ошибке сервера.
\end{example}

\subsection{Проверка непротиворечивости}

\begin{definition}
Спецификация называется \textbf{непротиворечивой}, если из неё нельзя вывести одновременно некоторое утверждение $P$ и его отрицание $\neg P$.
\end{definition}

Формально:
\[
    Spec \not\models P \land \neg P.
\]

\begin{example}
Противоречивая спецификация:
\begin{enumerate}
    \item При трёх неверных попытках входа пользователь блокируется.
    \item Пользователь никогда не блокируется при неверных попытках входа.
\end{enumerate}
\end{example}

\subsection{Проверка достижимости состояний}

Для автоматных спецификаций важно проверить, что все состояния достижимы из начального.

\begin{definition}
Состояние $q$ называется \textbf{достижимым}, если существует последовательность входов $w \in X^*$ такая, что:
\[
    \delta^*(q_0,w) = q.
\]
\end{definition}

\subsubsection{Алгоритм поиска достижимых состояний}

\textbf{Вход:}
\begin{itemize}
    \item множество состояний $Q$;
    \item начальное состояние $q_0$;
    \item отношение переходов $R \subseteq Q \times Q$.
\end{itemize}

\textbf{Выход:}
множество достижимых состояний $Reach$.

\begin{lstlisting}[style=pseudo]
Reach := {q0}
Queue := [q0]

while Queue не пуста do
    q := извлечь_из_Queue()

    for каждое состояние p такое, что (q,p) in R do
        if p not in Reach then
            добавить p в Reach
            добавить p в Queue
        end
    end
end

return Reach
\end{lstlisting}

После выполнения алгоритма состояния из множества:
\[
    Q \setminus Reach
\]
являются недостижимыми.

\begin{example}
Пусть:
\[
    Q = \{A,B,C,D\},
\]
\[
    q_0 = A,
\]
\[
    R = \{(A,B),(B,C)\}.
\]

Тогда:
\[
    Reach = \{A,B,C\}.
\]

Состояние $D$ недостижимо.
\end{example}

\subsection{Проверка детерминированности}

Автоматная спецификация детерминирована, если для каждой пары ``состояние--событие'' существует не более одного перехода.

Формально:
\[
    \forall q \in Q,\ \forall x \in X:
    |\{q' \mid \delta(q,x)=q'\}| \leq 1.
\]

\begin{example}
Недетерминированная таблица:

\[
\begin{array}{c|c|c}
\text{Состояние} & \text{Событие} & \text{Новое состояние} \\
\hline
A & x & B \\
A & x & C
\end{array}
\]

Для одной и той же пары $(A,x)$ указаны два разных результата.
\end{example}

\subsection{Проверка отсутствия тупиков}

\begin{definition}
\textbf{Тупиковое состояние} --- это состояние, из которого нет допустимых переходов, хотя система не должна завершаться.
\end{definition}

Формально состояние $q$ является тупиковым, если:
\[
    \neg Final(q) \land \nexists q' : (q,q') \in R.
\]

Проверка:
\begin{enumerate}
    \item построить граф переходов;
    \item найти все достижимые состояния;
    \item среди достижимых состояний найти такие, из которых нет исходящих переходов;
    \item исключить корректные финальные состояния;
    \item оставшиеся состояния считать дефектами спецификации.
\end{enumerate}

\subsection{Проверка инвариантов}

Инвариант $I$ должен выполняться:
\begin{enumerate}
    \item в начальном состоянии;
    \item сохраняться каждым переходом.
\end{enumerate}

Формально для переходной системы:
\[
    TS = (S, S_0, R),
\]
где:
\begin{itemize}
    \item $S$ --- множество состояний;
    \item $S_0 \subseteq S$ --- множество начальных состояний;
    \item $R \subseteq S \times S$ --- отношение переходов.
\end{itemize}

Инвариант $I$ корректен, если:
\[
    \forall s_0 \in S_0: I(s_0)
\]
и
\[
    \forall s,s' \in S: I(s) \land R(s,s') \Rightarrow I(s').
\]

\begin{theorem}[Индуктивная проверка инварианта]
Пусть для переходной системы $TS=(S,S_0,R)$ выполнены условия:
\[
    \forall s_0 \in S_0: I(s_0),
\]
\[
    \forall s,s' \in S: I(s) \land R(s,s') \Rightarrow I(s').
\]
Тогда инвариант $I$ выполняется во всех достижимых состояниях системы.
\end{theorem}

\begin{proof}
Докажем по индукции по длине пути из начального состояния.

\textbf{База индукции.}
Путь длины $0$ состоит только из начального состояния $s_0 \in S_0$.
По первому условию:
\[
    I(s_0)
\]
истинно.

\textbf{Индукционный переход.}
Пусть для любого состояния $s$, достижимого за $k$ шагов, истинно $I(s)$.
Рассмотрим состояние $s'$, достижимое за $k+1$ шаг.
Тогда существует состояние $s$, достижимое за $k$ шагов, такое что:
\[
    R(s,s').
\]

По предположению индукции:
\[
    I(s).
\]

По второму условию:
\[
    I(s) \land R(s,s') \Rightarrow I(s').
\]

Следовательно:
\[
    I(s')
\]
истинно.

По принципу математической индукции $I$ выполняется во всех достижимых состояниях.
\end{proof}

\section{Модельная проверка}

\begin{definition}
\textbf{Модельная проверка} --- это метод автоматической проверки конечной модели системы на соответствие формальной спецификации, обычно записанной во временной логике.
\end{definition}

Общая постановка:
\[
    M \models \varphi,
\]
где:
\begin{itemize}
    \item $M$ --- модель системы;
    \item $\varphi$ --- требуемое свойство;
    \item $\models$ --- отношение выполнения.
\end{itemize}

\subsection{Переходная система}

\begin{definition}
\textbf{Переходная система} задаётся кортежем:
\[
    M = (S, S_0, R, L),
\]
где:
\begin{itemize}
    \item $S$ --- множество состояний;
    \item $S_0 \subseteq S$ --- множество начальных состояний;
    \item $R \subseteq S \times S$ --- отношение переходов;
    \item $L: S \to 2^{AP}$ --- функция разметки атомарными высказываниями.
\end{itemize}
\end{definition}

\subsection{Примеры свойств}

\begin{itemize}
    \item \textbf{Безопасность}: плохое состояние никогда не достижимо.
    \[
        AG \neg error.
    \]

    \item \textbf{Живость}: если запрос возник, то когда-нибудь будет ответ.
    \[
        AG(request \Rightarrow AF response).
    \]

    \item \textbf{Достижимость}: существует путь к состоянию успеха.
    \[
        EF success.
    \]
\end{itemize}

\subsection{Алгоритм проверки свойства достижимости}

Пусть требуется проверить:
\[
    EF\ p,
\]
то есть ``существует путь, на котором когда-нибудь достижимо состояние, удовлетворяющее $p$''.

\textbf{Идея:}
надо проверить, достижимо ли из начального состояния хотя бы одно состояние, где истинно $p$.

\begin{lstlisting}[style=pseudo]
Visited := empty_set
Queue := S0

while Queue не пуста do
    s := извлечь_из_Queue()

    if p истинно в s then
        return true
    end

    if s not in Visited then
        добавить s в Visited

        for каждое s' такое, что R(s,s') do
            добавить s' в Queue
        end
    end
end

return false
\end{lstlisting}

\begin{example}
Пусть есть система:
\[
    S = \{s_0,s_1,s_2\},
\]
\[
    R = \{(s_0,s_1),(s_1,s_2)\}.
\]

Пусть:
\[
    p
\]
истинно только в $s_2$.

Тогда:
\[
    EF\ p
\]
истинно, потому что существует путь:
\[
    s_0 \to s_1 \to s_2.
\]
\end{example}

\subsection{Достоинства модельной проверки}

\begin{itemize}
    \item высокая степень автоматизации;
    \item возможность получить контрпример;
    \item хорошо подходит для протоколов и реактивных систем;
    \item позволяет обнаруживать тонкие ошибки в параллельных системах.
\end{itemize}

\subsection{Недостатки модельной проверки}

\begin{itemize}
    \item проблема взрыва числа состояний;
    \item необходимость построения формальной модели;
    \item трудность моделирования больших программ напрямую;
    \item зависимость результата от корректности модели.
\end{itemize}

\section{Доказательство теорем}

\begin{definition}
\textbf{Доказательство теорем} --- это метод проверки спецификаций и программ, при котором корректность выводится в формальной логической системе.
\end{definition}

Используются:
\begin{itemize}
    \item логика первого порядка;
    \item теория множеств;
    \item теория типов;
    \item логика Хоара;
    \item интерактивные доказатели: Coq, Isabelle/HOL, Lean;
    \item SMT-решатели: Z3, CVC5.
\end{itemize}

Проверка корректности сводится к доказательству утверждений вида:
\[
    Spec \Rightarrow Property.
\]

Или для программы:
\[
    \{P\}\ Program\ \{Q\}.
\]

\subsection{Общий алгоритм дедуктивной верификации}

\begin{enumerate}
    \item Записать спецификацию в формальном виде.
    \item Аннотировать программу предусловиями, постусловиями и инвариантами.
    \item Сгенерировать условия корректности.
    \item Передать условия корректности в SMT-решатель или доказатель теорем.
    \item Если условия доказаны, программа корректна относительно спецификации.
    \item Если доказательство не найдено, анализируется контрпример или усиливаются аннотации.
\end{enumerate}

\section{Анимация и исполнимые спецификации}

\begin{definition}
\textbf{Анимация спецификации} --- это выполнение или имитация формальной спецификации на конкретных данных с целью проверки её поведения.
\end{definition}

Анимация не доказывает корректность полностью, но помогает выявить:
\begin{itemize}
    \item неверно понятые требования;
    \item пропущенные случаи;
    \item ошибки в переходах;
    \item некорректные ограничения.
\end{itemize}

\begin{example}
Если спецификация банкомата задана автоматом, можно выполнить сценарий:
\[
    insertCard \to enterPIN \to chooseWithdraw \to takeCash.
\]

Если после события \texttt{takeCash} автомат не возвращает карту, это может указывать на ошибку спецификации.
\end{example}

\section{Тестирование на основе спецификации}

\begin{definition}
\textbf{Тестирование на основе спецификации} --- это построение тестов непосредственно из требований, модели или формальной спецификации.
\end{definition}

Примеры:
\begin{itemize}
    \item разбиение входной области на классы эквивалентности;
    \item анализ граничных значений;
    \item генерация тестов по автоматной модели;
    \item проверка контрактов во время выполнения.
\end{itemize}

\subsection{Генерация тестов по конечному автомату}

\textbf{Цель:}
получить набор последовательностей событий, покрывающих переходы автомата.

\begin{lstlisting}[style=pseudo]
Tests := empty_set

for каждый переход t = (q, event, q') do
    найти путь path от q0 до q
    test := path + [event]
    добавить test в Tests
end

return Tests
\end{lstlisting}

\begin{example}
Для турникета переход:
\[
    Locked \xrightarrow{coin} Unlocked
\]
покрывается тестом:
\[
    [coin].
\]

Переход:
\[
    Unlocked \xrightarrow{push} Locked
\]
покрывается тестом:
\[
    [coin, push].
\]
\end{example}

\section{Runtime verification}

\begin{definition}
\textbf{Runtime verification} --- это проверка поведения программы во время выполнения на соответствие формальной спецификации.
\end{definition}

В отличие от статической верификации, runtime verification анализирует конкретные трассы выполнения.

Спецификация может быть задана:
\begin{itemize}
    \item автоматом;
    \item контрактами;
    \item временной логикой;
    \item регулярными выражениями над событиями.
\end{itemize}

Пример свойства:
\[
    open \Rightarrow eventually\ close.
\]

То есть если ресурс был открыт, он должен быть когда-нибудь закрыт.

\section{Сравнение методов спецификации}

\[
\begin{array}{|p{4cm}|p{5cm}|p{5cm}|}
\hline
\textbf{Метод} & \textbf{Достоинства} & \textbf{Недостатки} \\
\hline
Естественный язык & Простота, понятность & Неоднозначность, трудность проверки \\
\hline
Блок-схемы & Наглядность потока управления & Плохо масштабируются \\
\hline
Структурные программы & Читаемость, иерархичность & Не всегда удобны для событийных систем \\
\hline
Автоматы & Хороши для состояний и событий & Возможен взрыв числа состояний \\
\hline
Логика Хоара & Строгие доказательства корректности & Требует инвариантов и математической подготовки \\
\hline
Временные логики & Удобны для реактивных систем & Сложность формулировки свойств \\
\hline
Z, VDM, B & Строгая математическая спецификация & Высокий порог входа \\
\hline
UML & Удобство коммуникации & Не всегда строгая семантика \\
\hline
\end{array}
\]

\section{Сравнение методов проверки спецификаций}

\[
\begin{array}{|p{4cm}|p{5cm}|p{5cm}|}
\hline
\textbf{Метод проверки} & \textbf{Что обнаруживает} & \textbf{Ограничения} \\
\hline
Инспекция & Неясности, противоречия, пропуски & Зависит от экспертов \\
\hline
Проверка типов & Ошибки согласования данных & Не проверяет смысл требований \\
\hline
Анимация & Ошибки сценариев & Не даёт полного доказательства \\
\hline
Тестирование по спецификации & Несоответствие реализации требованиям & Покрывает только выбранные случаи \\
\hline
Модельная проверка & Нарушения свойств в конечной модели & Взрыв числа состояний \\
\hline
Доказательство теорем & Строгая корректность & Трудоёмкость \\
\hline
Runtime verification & Ошибки в реальном выполнении & Проверяет только наблюдаемые трассы \\
\hline
\end{array}
\]

\section{Итоговая схема связи понятий}

\begin{center}
\begin{tikzpicture}[
    node distance=1.6cm,
    box/.style={draw, rounded corners, align=center, minimum width=4cm, minimum height=0.9cm},
    >=Latex
]
    \node[box] (req) {Требования};
    \node[box, below of=req] (spec) {Спецификация};
    \node[box, below left of=spec, xshift=-2.5cm] (formal) {Формальная\\спецификация};
    \node[box, below right of=spec, xshift=2.5cm] (semi) {Полуформальная\\спецификация};
    \node[box, below of=formal] (verify) {Верификация};
    \node[box, below of=semi] (visual) {Диаграммы,\\автоматы, схемы};
    \node[box, below of=verify] (impl) {Реализация};
    \node[box, below of=visual] (test) {Тестирование\\и анимация};

    \draw[->] (req) -- (spec);
    \draw[->] (spec) -- (formal);
    \draw[->] (spec) -- (semi);
    \draw[->] (formal) -- (verify);
    \draw[->] (semi) -- (visual);
    \draw[->] (verify) -- (impl);
    \draw[->] (visual) -- (test);
    \draw[->] (test) -- (impl);
\end{tikzpicture}
\end{center}

\section{Краткие выводы}

\begin{enumerate}
    \item Спецификация фиксирует требуемое поведение программы.
    \item Неформальные спецификации удобны, но неоднозначны.
    \item Формальные спецификации позволяют проводить строгую проверку.
    \item Схемное программирование задаёт алгоритм как граф управления.
    \item Структурное программирование ограничивает управление последовательностью, ветвлением и циклом.
    \item Теорема Бёма--Якопини обосновывает достаточность этих трёх конструкций.
    \item Визуальное программирование использует графическое представление программ и моделей.
    \item Автоматное программирование описывает поведение через состояния, события, переходы и действия.
    \item Абстрактный автомат является математической моделью, а управляющий автомат выполняет функцию управления операционной частью или объектом.
    \item Проверка спецификаций включает инспекцию, анализ полноты, непротиворечивости, достижимости, модельную проверку, доказательство теорем, анимацию и тестирование.
\end{enumerate}

\end{document}
